\Article{341}
Primum caput harum investigationum in eo consistet, ut functionum primarum cuiusvis dimensionis multitudinem determinemus. Quoniam enim pro modulo determinato numerus omnium functionum diversarum (incongruarum) cuiuslibet gradus est definitus, ex his vero aliae sunt ex primis inferiorum graduum compositae, aliae primae, etiam harum numerus finitus erit. Rigorosa huius rei evolutio satis est lubrica; a casibus simplicioribus incipiemus.

Posito modulo $=p$, numerus omnium functionum diversarum $n^{\mathrm{ti}}$ gradus huius formae
\[
x^n+Ax^{n-1}+Bx^{n-2}+Cx^{n-3}+\text{etc.}
\]
erit $p^n$; coëfficientium enim $\pdot A,B\pdot C$ etc. numerus est $n$; et quum quivis independenter a reliquis possit esse $\equiv0,1,2,3\dotstwo(p-1)(\mathrm{mod.}\,p)$, ex combinationum theoria sequitur, $p^n$ combinationes diversas haberi, quae igitur omnium functionum diversarum huius gradus complexum definiunt.

Ita functiones unius dimensionis erunt $p$, scilicet $x,x+1,x+2$ usque ad $x+p-1$; functiones duarum dimensionum $pp$ etc.

\Article{342}
Iam supra monuimus, omnes functiones primi gradus pro primis habendas esse; si igitur, quod ad propositum nostrum sufficit, ad eas functiones nos restringamus, quarum terminus summus habet coëfficientem $1$, erunt $p$ functiones primi gradus seu unius dimensionis.

Functiones secundi gradus omnes aut e binis primi gradus erunt compositae aut primae. Iam ex combinationum theoria constat, $p$ res diversas admissis repetitionibus $\frac{p\pdot p+1}{1\pdot2}$ modis diversis combinari posse, quare totidem functiones erunt e binis primis unius dimensionis compositae, adeoque $pp-\frac{p\pdot p+1}{1\pdot2}=\frac12(pp-p)$ functiones primae duarum dimensionum.

\SourceLeaf{230}
Simili modo e functionibus omnibus tertii gradus, quarum numerus est $p^3$, excludendae sunt eae, quae e ternis primis unius dimensionis componuntur, quarum numerus est $\frac{p\pdot p+1\pdot p+2}{1\pdot2\pdot3}$; insuperque eae, quae e functione prima unius aliaque duarum dimensionum componuntur, quarum numerus est $p\pdot\frac12(pp-p)$; quibus deletis restabunt $\frac13(p^3-p)$; tot igitur sunt primae trium dimensionum. Elucet hoc modo semper continuari posse.

\Article{343}
Ut autem hae operationes facilius absolvantur simulque ad evolutionem legis generalis via sternatur, rem generaliter considerabimus. Brevitatis gratia designamus per $(1)$ multitudinem functionum primarum unius dimensionis, per $(2)$ numerum functionum primarum duarum dimensionum, sic porro per $(1^2)$ multitudinem functionum e binis primis unius dimensionis compositarum etc. etc., generaliter per $(1^\alpha2^\beta3^\gamma\ldots)$ multitudinem functionum omnium, quae e functionibus primis compositae sunt, scilicet ex $\alpha$ unius, $\beta$ duarum, $\gamma$ trium etc. dimensionum, quarum itaque dimensio erit $\alpha+2\beta+3\gamma+\text{etc.}$ Tum per praecedentia theoriamque combinationum elucet, fore
\[
(1^\alpha2^\beta3^\gamma4^\delta\ldots)=(1^\alpha)(2^\beta)(3^\gamma)(4^\delta)\dotstwo
\]
\[
(1^\alpha)=\frac{(1)\pdot(1)+1\pdot(1)+2\pdot(1)+3\ldots(1)+\alpha-1}{1\pdot2\pdot3\pdot4\ldots\alpha}
\]
seu generaliter
\[
(a^\alpha)=\frac{(a)\pdot(a)+1\pdot(a)+2\pdot(a)+3\ldots(a)+\alpha-1}{1\pdot2\pdot3\pdot4\ldots\alpha}
\]

Denique manifestum est, si omnes modi diversi numerum $n$ e numeris $1,2,3,\ldots$ per additionem componendi colligantur, qui designentur per $\alpha\pdot1+\beta\pdot2+\gamma\pdot3+\text{etc.}$, summam omnium harum expressionum $(1^\alpha2^\beta3^\gamma\dotstwo)$ aequalem fore multitudini omnium functionum $n$ dimensionum, i. e. $=p^n$. Ita
\[
\begin{aligned}
p&=(1)\\
pp&=(1^2)+(2)\\
p^3&=(1^3)+(1\pdot2)+(3)\\
p^4&=(1^4)+(1^2\pdot2)+(1\pdot3)+(2^2)+(4)\\
&\text{etc.}
\end{aligned}
\]

Perspicuum est, in expressione $p^n$ praeter quantitates $(1),(2),(3)$ etc. etiam hanc
\SourceLeaf{231}
ingredi $(n)$, unde patet, quomodo omnes quantitates per praecedentes sint determinandae. Ita invenitur
\[
\begin{array}{lll}
(1)=p&(4)=\frac14(p^4-pp)&(7)=\frac17(p^7-p)\\
(2)=\frac12(pp-p)&(5)=\frac15(p^5-p)&(8)=\frac18(p^8-p^4)\\
(3)=\frac13(p^3-p)&(6)=\frac16(p^6-p^3-pp+p)&\text{etc.}
\end{array}
\]

\Article{344--346}
Observatur ex hoc seriei initio, summum terminum expressionis $(n)$ esse $\frac1n p^n$, ad quem, si $n$ est primus, accedit $-\frac1n p$; at si $n$ est compositus, lex minus elucet. Si vero attentius rem consideramus, videmus esse
\[
\begin{array}{ll}
p=(1)&p^5=5(5)+(1)\\
pp=2(2)+(1)&p^6=6(6)+3(3)+2(2)+(1)\\
p^3=3(3)+(1)&p^7=7(7)+(1)\\
p^4=4(4)+2(2)+(1)&p^8=8(8)+4(4)+2(2)+(1)\text{ etc.}
\end{array}
\]
ubi lex progressionis est manifesta; scilicet si omnes numeri $n$ divisores sint $\alpha,\beta,\gamma,\delta$ etc., erit
\[
p^n=\alpha(\alpha)+\beta(\beta)+\gamma(\gamma)+\delta(\delta)+\text{etc.}
\]

Huius observationis generalitatem iam demonstrare accingimur.

Ostendimus summam omnium talium expressionum $(1^\alpha)(2^\beta)(3^\gamma)\ldots$ si semper $\alpha+2\beta+3\gamma+\ldots=n$, exhaurire omnes functiones $n$ dimensionum adeoque esse $=p^n$. Hinc patet, \textemdash\ \textemdash\ \textemdash. Si
\[
\left(\frac1{1-x}\right)^{(1)}\left(\frac1{1-xx}\right)^{(2)}\left(\frac1{1-x^3}\right)^{(3)}\ldots
\]
evolvatur in seriem $1+Ax+Bx^2\ldots=P$, erit
\[
A=p,\quad B=p^2,\quad C=p^3\text{ etc.}
\]
\[
\frac{x\,\mathrm{d}P}{P\,\mathrm{d}x}=\frac{(1)x}{1-x}+\frac{2(2)x^3}{1-x^2}+\frac{3(3)x^3}{1-x^3}\ldots
\]
\noindent\hspace*{.28\linewidth}\textemdash\ \textemdash\ \textemdash\ \textemdash

[hinc substituendo $\frac{px}{1-px}$ pro $\frac{x\,\mathrm{d}P}{P\,\mathrm{d}x}$ et evolvendo singulas fractiones in series infinitas theorematis veritas sponte elucet.]

\SourceLeaf{232}
\Article{347}
Theorema hoc etiam alio modo exprimi potest. Scilicet si numeri $n$ divisores omnes sint $n,1,\delta,\delta',\delta'',\delta'''$ etc., theorema in eo consistit, ut sit
\[
p^n=n(n)+(1)+\delta(\delta)+\delta'(\delta')+\text{etc.}
\]
Iam patet, productum ex $(n)$ functionibus primis, quae sunt $n$ dimensionum, habere $n(n)$ dimensiones et sic de reliquis, quare

\emph{Productum ex omnibus functionibus primis dimensionis unius, dimensionum $n,\delta,\delta'$ etc. habebit $p^n$ dimensiones.}

Facile nunc est ex hoc theoremate valorem expressionis $(n)$ ipsam deducere, sed brevitatis gratia analysin, quae non est difficilis, supprimimus. Sit itaque $n=a^\alpha b^\beta c^\gamma$ etc., ita ut $a,b,c$ etc. sint numeri primi diversi, eritque
\[
n(n)=p^n-\Sigma p^{\frac na}+\Sigma p^{\frac n{ab}}-\Sigma p^{\frac n{abc}}\text{ etc.}
\]
ubi $\Sigma p^{\frac n{abc\dotstwo}}$ significat complexum omnium expressionum huic $p^{\frac n{abc\dotstwo}}$ similium, si quantitates $a,b,c\dotstwo$ quomodocunque inter se permutentur. Ita pro $n=36$ erit $36(36)=p^{36}-p^{18}-p^{12}+p^6$.

Unam adhuc observationem adiicere liceat. Si $n$ est formae $a^\alpha$ et $a$ primus, erit $n(n)=p^n-p^{\frac na}$, quare, quum $(n)$ necessario sit integer, erit quicquid sit $p$,
\[
p^n\equiv p^{\frac na}\;(\mathrm{mod.}\,n)
\]
quare, si $p$ ad $a$ primus erit,
\[
p^{n-\frac na}\equiv1\;(\mathrm{mod.}\,n)
\]
et pro $\alpha=1$
\[
p^{a-1}\equiv1\;(\mathrm{mod.}\,a)
\]
Memorabile est, haec theoremata tam diversis modis erui posse.
